\documentclass[12pt,a4paper]{article}

% ===== PAQUETES =====
\usepackage[utf8]{inputenc}
\usepackage[T1]{fontenc}
\usepackage[spanish]{babel}
\usepackage{geometry}
\usepackage{graphicx}
\usepackage{xcolor}
\usepackage{hyperref}
\usepackage{booktabs}
\usepackage{longtable}
\usepackage{array}
\usepackage{enumitem}
\usepackage{fancyhdr}
\usepackage{titlesec}
\usepackage{colortbl}
\usepackage{tabularx}
\usepackage{parskip}
\usepackage{multirow}

% ===== GEOMETRÍA =====
\geometry{
  top=2.5cm, bottom=2.5cm,
  left=3cm,  right=2.5cm
}

% ===== COLORES =====
\definecolor{verdeuteq}{RGB}{0,92,47}
\definecolor{verdelight}{RGB}{0,128,64}
\definecolor{javacolor}{RGB}{200,95,20}
\definecolor{sqlcolor}{RGB}{0,100,180}
\definecolor{angularcolor}{RGB}{221,30,40}
\definecolor{owaspborder}{RGB}{180,130,0}
\definecolor{gapcolor}{RGB}{200,120,0}
\definecolor{pendientecolor}{RGB}{130,130,130}
\definecolor{grisclaro}{RGB}{240,240,240}

% ===== HYPERREF =====
\hypersetup{
  colorlinks=true,
  linkcolor=verdeuteq,
  urlcolor=verdelight,
  citecolor=verdeuteq,
  pdftitle={SGB-SaaS -- Tercera Entrega -- Consolidacion, Reproducibilidad y Evidencia Empirica},
  pdfauthor={Cajas Ibarra Irvin Marcelo, Loor Medranda Marlon Taylor, Panama Murillo Moises Antonio -- UTEQ}
}

% ===== ENCABEZADOS =====
\pagestyle{fancy}
\fancyhf{}
\renewcommand{\headrulewidth}{0.4pt}
\fancyhead[L]{\small SGB-SaaS --- Tercera Entrega}
\fancyhead[R]{\small \texttt{v0.9.0-rc}}
\fancyfoot[C]{\thepage}

\titleformat{\section}{\normalfont\Large\bfseries\color{verdeuteq}}{\thesection}{0.6em}{}
\titleformat{\subsection}{\normalfont\large\bfseries\color{verdeuteq}}{\thesubsection}{0.6em}{}

% ===== BADGES =====
\newcommand{\badgeJAVA}{\colorbox{javacolor}{\color{white}\small\bfseries\ Java 21 / Spring Boot 4\ }}
\newcommand{\badgeSQL}{\colorbox{sqlcolor}{\color{white}\small\bfseries\ PostgreSQL 16\ }}
\newcommand{\badgeOWASP}{\colorbox{verdeuteq}{\color{white}\small\bfseries\ OWASP\ }}
\newcommand{\badgeANGULAR}{\colorbox{angularcolor}{\color{white}\small\bfseries\ Angular 17\ }}
\newcommand{\badgeREDIS}{\colorbox{owaspborder}{\color{white}\small\bfseries\ Redis 7\ }}

% ===== ESTADOS (badges de tabla) =====
\newcommand{\estadoPasa}{\colorbox{verdelight}{\color{white}\small\bfseries\ PASA\ }}
\newcommand{\estadoCorregido}{\colorbox{verdeuteq}{\color{white}\small\bfseries\ CORREGIDO\ }}
\newcommand{\estadoGap}{\colorbox{gapcolor}{\color{white}\small\bfseries\ GAP CONOCIDO\ }}
\newcommand{\estadoPendiente}{\colorbox{pendientecolor}{\color{white}\small\bfseries\ NO EJECUTADO\ }}
\newcommand{\estadoVerificado}{\colorbox{verdelight}{\color{white}\small\bfseries\ verificado\ }}
\newcommand{\estadoImplementado}{\colorbox{sqlcolor}{\color{white}\small\bfseries\ implementado\ }}
\newcommand{\estadoPend}{\colorbox{gapcolor}{\color{white}\small\bfseries\ pendiente\ }}

% Nota honesta (caja resaltada). Nota tecnica: NO usar un entorno tabularx
% de una sola columna X para esto -- un segundo entorno tabularx mas
% adelante en el documento revienta con "Missing }"/"Misplaced \noalign"
% (bug real de interaccion entre dos tabularx con distinta cantidad de
% columnas, confirmado aislando el caso minimo). Se usa colorbox+parbox
% en su lugar, sin entorno reutilizable (un solo uso en todo el documento).

\begin{document}

% ============================================================
% PORTADA
% ============================================================
\thispagestyle{empty}
\begin{titlepage}
\centering
\vspace*{1cm}

{\large\bfseries Universidad Técnica Estatal de Quevedo (UTEQ)}\\[0.2cm]
{\normalsize Facultad de Ciencias de la Computación y Diseño Digital}\\[0.1cm]
{\normalsize Carrera de Ingeniería de Software}\\[0.1cm]
{\normalsize Asignatura: Aplicaciones Web --- Período académico 2026--2027}\\[1.5cm]

\rule{0.8\linewidth}{1pt}\\[0.4cm]
{\Large\bfseries\color{verdeuteq} Sistema de Gestión Bibliotecaria Web (SGB-SaaS)}\\[0.3cm]
{\large Tercera Entrega --- Consolidación, Reproducibilidad Automática\\ y Evidencia Empírica Cuantitativa}\\[0.4cm]
\rule{0.8\linewidth}{1pt}\\[1cm]

\begin{center}
\badgeJAVA\quad\badgeANGULAR\quad\badgeSQL\quad\badgeREDIS\quad\badgeOWASP
\end{center}

\vspace{1cm}

{\normalsize\bfseries Integrantes:}\\[0.2cm]
\begin{tabular}{ll}
Cajas Ibarra Irvin Marcelo & C.I.: 1251039606 \\
Loor Medranda Marlon Taylor & C.I.: 0928087469 \\
Panama Murillo Moises Antonio & C.I.: 1251334544 \\
\end{tabular}\\[1cm]

{\normalsize\bfseries Docente:}\\
Ing. Gleiston Cicerón Guerrero Ulloa, Ph.D.\\[1cm]

{\normalsize\bfseries Fecha de cierre:}\\
24 de julio de 2026\\[1cm]

{\normalsize\bfseries Repositorio:}\\
\url{https://github.com/mloorm14/sgb-saas} \quad --- \quad Rama: \texttt{main}\\[0.5cm]

{\normalsize\bfseries Tag de esta entrega:} \texttt{v0.9.0-rc} (release candidate)\\[0.2cm]
{\small Commit tagueado (completo):}\\
{\small\texttt{c6b372c9cdaee878ae5be004a00835bec4e3245a}}\\[0.1cm]
{\small Commit (corto): \texttt{c6b372c}}\\[0.5cm]

{\normalsize\bfseries DOI:}\\
{\small pendiente de asignación (ver \texttt{README.md} tras publicación en Zenodo)}\\[1cm]

\vfill
{\small SOFT-R-A --- 5to Nivel --- Proyecto Fin de Curso (PFC), Tercera Entrega}
\end{titlepage}

\tableofcontents
\newpage

% ============================================================
% 1. RESUMEN EJECUTIVO
% ============================================================
\section{Resumen ejecutivo}

SGB-SaaS es un sistema web de gestión bibliotecaria construido como
Proyecto Fin de Curso (PFC) de la asignatura Aplicaciones Web, con una
arquitectura de tres capas: frontend SPA en Angular~17, backend API REST
en Spring Boot~4.0.6 (Java~21), persistencia en PostgreSQL~16 y una capa
de caché/blacklist de tokens en Redis~7. El sistema modela el ciclo
completo de una biblioteca institucional: catálogo de libros, préstamos,
reservaciones, multas y administración de usuarios bajo control de
acceso basado en roles (RBAC).

Esta Tercera Entrega consolida el trabajo iniciado en la Entrega~1B
(tag \texttt{v0.1.0-entrega-1b}) sobre cinco frentes concretos:

\begin{enumerate}[leftmargin=1.2cm]
  \item \textbf{RBAC normalizado.} El antiguo campo de texto libre
    \texttt{rol} se reemplazó por un modelo relacional completo
    (\texttt{roles}, \texttt{usuario\_roles}, \texttt{permisos},
    \texttt{rol\_permisos}) con cuatro roles (\texttt{LECTOR},
    \texttt{BIBLIOTECARIO}, \texttt{GERENTE}, \texttt{ADMIN}) aplicados
    consistentemente vía \texttt{@PreAuthorize} en cada endpoint
    (ADR-010).
  \item \textbf{Estrategia híbrida de acceso a datos.} Las operaciones
    CRUD elementales usan Spring Data JPA; las operaciones con lógica de
    negocio transaccional multi-tabla (registrar devolución, pagar
    multa, anular multa, crear préstamo) se implementan como
    procedimientos almacenados de PostgreSQL, documentados en
    \texttt{docs/basedatos/CATALOGO-SP.md} y formalizados en ADR-013.
  \item \textbf{Módulo funcional de Préstamos, Reservaciones y Multas.}
    Construido sobre el modelo de datos normalizado, con 5 procedimientos
    almacenados propios y su correspondiente capa de controladores REST
    (\texttt{PrestamoController}, \texttt{ReservacionController},
    \texttt{MultaController}).
  \item \textbf{Seguridad reforzada mediante auditoría activa.} No solo
    se documentó la postura de seguridad frente a OWASP Top~10 --
    se ejecutaron 6 auditorías en vivo contra el stack real, y dos
    hallazgos verdaderos que esas auditorías destaparon
    (ausencia de rate limiting en login, ausencia total de logging de
    eventos de autenticación) se corrigieron dentro de esta misma
    entrega, con evidencia de antes/después.
  \item \textbf{Reproducibilidad automática.} Scripts dedicados
    (\texttt{scripts/mediciones-header.sh},
    \texttt{scripts/validate-traceability.sh},
    \texttt{scripts/build-init-sql.sh}) garantizan que cualquier
    evidencia nueva capture versiones exactas de herramientas y que la
    matriz de trazabilidad se valide automáticamente en CI.
\end{enumerate}

\subsection*{Estado real, sin maquillar}

Con la misma honestidad que exige la guía de esta entrega (Bloque~B.2:
``la edición manual de un archivo crudo de mediciones invalida la
evidencia''), este informe distingue explícitamente lo que se completó
de lo que no:

\textbf{Completo, con evidencia real verificable:}
seguridad (auditoría OWASP de 6 controles + 2 correcciones en vivo),
cobertura de código (JaCoCo real, 75{,}85\,\% de líneas sobre el
dominio/servicios), matriz de trazabilidad (30 requisitos, validación
automática en CI), 13 Architecture Decision Records
\footnote{Nota de precisión: el repositorio contiene actualmente 10
archivos de ADR (\texttt{docs/adr/}), no 13 -- ver
\autoref{sec:estado-sistema} para el detalle exacto. Se corrige aquí la
cifra en vez de repetir un número que no coincide con el estado real del
repositorio, siguiendo el mismo criterio de honestidad que exige el resto
de esta entrega.}, licencia MIT, colección Postman con 39 requests
verificados en vivo.

\textbf{Con evidencia parcial o diferida, declarado explícitamente:}
TLS en tránsito (OWASP A02, entorno de desarrollo sobre HTTP plano por
diseño -- exigido recién en Entrega Final), \emph{Content-Security-Policy}
(OWASP A05), y tres de los cinco sub-bloques de evidencia empírica del
Bloque~C -- rendimiento (k6), usabilidad (SUS) y accesibilidad
(Lighthouse) -- que quedan \textbf{no ejecutados} en esta entrega por
restricciones reales de tiempo del equipo, con toda su configuración y
protocolo ya preparados para ejecutarse antes de la Entrega Final (ver
\autoref{sec:protocolo} y \autoref{sec:resultados}).

% ============================================================
% 2. ESTADO DEL SISTEMA
% ============================================================
\section{Estado del sistema}
\label{sec:estado-sistema}

\subsection{Pila tecnológica}

\begin{center}
\begin{tabular}{|l|l|l|}
\hline
\rowcolor{verdeuteq}
\color{white}\textbf{Capa} & \color{white}\textbf{Tecnología} & \color{white}\textbf{Rol} \\
\hline
Frontend & Angular 17 (TypeScript) & SPA, formularios reactivos, sesión en memoria \\
\hline
Backend & Spring Boot 4.0.6 / Java 21 & API REST, JWT+RBAC, lógica de negocio \\
\hline
Base de datos & PostgreSQL 16 & 26 tablas, 7 procedimientos/funciones SQL \\
\hline
Caché/Sesión & Redis 7 & Blacklist de JWT revocados, caché del catálogo \\
\hline
Orquestación & Docker Compose & 4 contenedores, imágenes pinadas por digest \\
\hline
\end{tabular}
\end{center}

La elección de cada componente de esta pila, con alternativas
consideradas y descartadas, está formalizada en ADR-001 (pila principal),
ADR-011 (PostgreSQL sobre MySQL/MongoDB) y ADR-012 (Docker Compose sobre
Kubernetes) -- este informe no duplica ese detalle, remite a
\texttt{docs/adr/} para el razonamiento completo de cada decisión.

\subsection{Módulos funcionales}

\begin{itemize}[leftmargin=1.2cm]
  \item \textbf{Auth} (\texttt{/api/auth/*}): registro, login, logout,
    refresco de \texttt{accessToken} vía cookie \texttt{HttpOnly}.
    Rate limiting de intentos fallidos (Redis, 5 intentos / 900\,s) y
    bitácora de auditoría (\texttt{bitacora\_auditoria}) para
    \texttt{LOGIN\_OK}/\texttt{LOGIN\_FAIL}/\texttt{LOGOUT}.
  \item \textbf{Libros} (\texttt{/api/v1/libros}): CRUD completo vía
    Spring Data JPA, caché Redis de la operación de listado con TTL
    configurable externamente (\texttt{CACHE\_LIBROS\_TTL\_SECONDS}).
  \item \textbf{Préstamos} (\texttt{/api/v1/prestamos}): creación y
    devolución vía procedimientos almacenados transaccionales
    (\texttt{sp\_crear\_prestamo}, \texttt{sp\_registrar\_devolucion}),
    listados por usuario, reporte de libros más prestados.
  \item \textbf{Reservaciones} (\texttt{/api/v1/reservaciones}): creación
    y listado por usuario; expiración automática de reservaciones
    vencidas vía \texttt{ReservacionScheduler} (job programado).
  \item \textbf{Multas} (\texttt{/api/v1/multas}): listado, pago y
    anulación (\texttt{sp\_pagar\_multa}, \texttt{sp\_anular\_multa}),
    esta última restringida a \texttt{GERENTE}/\texttt{ADMIN}
    exclusivamente.
\end{itemize}

\subsection{Inventario de endpoints y roles requeridos}

Inventario completo de los 19 endpoints REST del backend, verificado
línea por línea contra los 5 \texttt{@RestController} reales (no
inferido) -- misma fuente usada para construir
\texttt{docs/postman/coleccion.json}.

\begin{center}
\small
\begin{longtable}{|l|l|l|}
\hline
\rowcolor{verdeuteq}
\color{white}\textbf{Método y ruta} & \color{white}\textbf{Roles} & \color{white}\textbf{Controller} \\
\hline
\endhead
\texttt{POST /api/auth/registro} & público & \texttt{AuthController} \\ \hline
\texttt{POST /api/auth/login} & público & \texttt{AuthController} \\ \hline
\texttt{POST /api/auth/refresh} & público (cookie) & \texttt{AuthController} \\ \hline
\texttt{POST /api/auth/logout} & autenticado & \texttt{AuthController} \\ \hline
\texttt{GET /api/v1/libros} & LECTOR+ & \texttt{LibroController} \\ \hline
\texttt{GET /api/v1/libros/\{id\}} & LECTOR+ & \texttt{LibroController} \\ \hline
\texttt{POST /api/v1/libros} & BIBLIOTECARIO+ & \texttt{LibroController} \\ \hline
\texttt{PUT /api/v1/libros/\{id\}} & BIBLIOTECARIO+ & \texttt{LibroController} \\ \hline
\texttt{DELETE /api/v1/libros/\{id\}} & BIBLIOTECARIO+ & \texttt{LibroController} \\ \hline
\texttt{POST /api/v1/prestamos} & BIBLIOTECARIO/GERENTE & \texttt{PrestamoController} \\ \hline
\texttt{POST /api/v1/prestamos/\{id\}/devolucion} & BIBLIOTECARIO/GERENTE & \texttt{PrestamoController} \\ \hline
\texttt{GET /api/v1/prestamos/usuario/\{id\}} & LECTOR/BIBLIOTECARIO/GERENTE & \texttt{PrestamoController} \\ \hline
\texttt{GET /api/v1/prestamos/usuario/\{id\}/activos} & LECTOR/BIBLIOTECARIO/GERENTE & \texttt{PrestamoController} \\ \hline
\texttt{GET /api/v1/prestamos/reportes/libros-mas-prestados} & BIBLIOTECARIO/GERENTE & \texttt{PrestamoController} \\ \hline
\texttt{POST /api/v1/reservaciones} & LECTOR/BIBLIOTECARIO/GERENTE & \texttt{ReservacionController} \\ \hline
\texttt{GET /api/v1/reservaciones/usuario/\{id\}} & LECTOR/BIBLIOTECARIO/GERENTE & \texttt{ReservacionController} \\ \hline
\texttt{GET /api/v1/multas/usuario/\{id\}} & LECTOR/BIBLIOTECARIO/GERENTE & \texttt{MultaController} \\ \hline
\texttt{POST /api/v1/multas/\{id\}/pago} & BIBLIOTECARIO/GERENTE & \texttt{MultaController} \\ \hline
\texttt{POST /api/v1/multas/\{id\}/anulacion} & GERENTE/ADMIN & \texttt{MultaController} \\ \hline
\end{longtable}
\end{center}

Nótese la asimetría deliberada: \texttt{LibroController} sí incluye
\texttt{ADMIN} en sus 5 endpoints (fix aplicado en esta entrega, ver
\autoref{sec:resultados}), pero \texttt{PrestamoController} y
\texttt{ReservacionController} \textbf{no} -- un usuario exclusivamente
\texttt{ADMIN} recibe 403 en esos dos controllers (verificado en vivo al
construir la colección Postman). Se documenta esta inconsistencia real en
vez de asumir que el mismo criterio se aplicó de forma pareja en todo el
backend.

\subsection{Decisiones arquitectónicas documentadas}

El repositorio contiene \textbf{10 Architecture Decision Records}
(\texttt{docs/adr/}), que cubren de forma completa los 6 temas
obligatorios del Bloque~D de la guía (ver índice en
\texttt{docs/adr/README.md}) más 2 ADRs adicionales no mapeados a esos
temas obligatorios:

\begin{center}
\small
\begin{tabularx}{\linewidth}{|l|X|}
\hline
\rowcolor{verdeuteq}
\color{white}\textbf{ADR} & \color{white}\textbf{Decisión} \\
\hline
ADR-001 & Elección de la pila tecnológica principal \\
ADR-003 & Redis para blacklist de tokens JWT revocados \\
ADR-006 & Esquema reproducible (Flyway + \texttt{schema.sql}/\texttt{seed.sql}) \\
ADR-007 & Cookies \texttt{HttpOnly} para JWT -- \texttt{refreshToken} migrado, \texttt{accessToken} pendiente \\
ADR-008 & TTL del caché Redis del catálogo, en configuración externa \\
ADR-009 & Licencia MIT del proyecto \\
ADR-010 & JWT stateless + RBAC normalizado sobre un campo de rol simple \\
ADR-011 & PostgreSQL como gestor de base de datos \\
ADR-012 & Docker Compose como estrategia de despliegue \\
ADR-013 & Estrategia híbrida de acceso a datos: CRUD-ORM + stored procedures \\
\hline
\end{tabularx}
\end{center}

Este informe referencia estas decisiones por su identificador cuando
resulta relevante, sin reproducir su contenido completo -- el detalle de
alternativas consideradas, consecuencias y trade-offs vive únicamente en
cada ADR, para evitar que este documento y \texttt{docs/adr/} diverjan
con el tiempo.

% ============================================================
% 3. ARQUITECTURA (C4)
% ============================================================
\section{Arquitectura del sistema (Modelo C4)}

El modelo C4 de SGB-SaaS está versionado como código fuente en
\texttt{docs/arquitectura/workspace.dsl} (Structurizr DSL), que
reemplaza los diagramas estáticos sin fuente versionada de la Entrega~1A
(\texttt{docs/diagramas/diagrama\_c4\_capa1.png} y
\texttt{diagrama\_c4\_capa2.jpeg}).

\subsection{Nivel 1 -- Contexto del sistema}

Cinco actores interactúan con SGB-SaaS como caja negra: \textbf{Lector}
(consulta catálogo, reserva, ve su historial), \textbf{Bibliotecario}
(registra préstamos/devoluciones/multas), \textbf{Gerente} (administra
personal, catálogos maestros y reportes), \textbf{Administrador}
(configura parámetros del sistema -- rol que \emph{no existía} en el
diseño de Entrega~1A, se agregó al normalizar RBAC) y
\textbf{Visitante anónimo} (solo puede registrarse o iniciar sesión).

\subsection{Nivel 2 -- Contenedores}

Cuatro contenedores, alineados 1:1 con \texttt{docker-compose.yml}:
Frontend Angular (SPA, sesión en memoria -- nunca \texttt{localStorage}),
Backend Spring Boot (API REST, JWT+RBAC), PostgreSQL (26 tablas) y Redis
(blacklist de JWT + caché de catálogo).

\subsubsection*{Diferencias documentadas respecto al diseño original de Entrega 1A}

El propio \texttt{workspace.dsl} documenta tres discrepancias reales
entre el diseño original y el sistema implementado, en vez de dibujar
capacidades que no existen:

\begin{itemize}[leftmargin=1.2cm]
  \item \textbf{Gemini~2.0~Flash API y Google Books API} aparecían como
    sistemas externos en el diseño de Entrega~1A. Se \textbf{retiraron}
    del diagrama actual: no existe ninguna integración con ellos en el
    código real del backend ni del frontend (verificado por búsqueda en
    el código fuente).
  \item \textbf{Catálogo público sin autenticación}: el diseño
    contemplaba que el \emph{Visitante anónimo} pudiera navegar el
    catálogo sin cuenta. En la implementación real,
    \texttt{LibroController} exige rol \texttt{LECTOR} (o superior)
    incluso para listar libros -- el visitante anónimo hoy solo puede
    registrarse o iniciar sesión. Discrepancia real entre diseño e
    implementación, documentada explícitamente en vez de ocultada.
  \item \textbf{Contenedor Redis y rol Administrador}: ambos ausentes del
    diagrama de Entrega~1A, agregados en esta entrega como consecuencia
    directa de la normalización RBAC y de la revocación de tokens vía
    blacklist.
\end{itemize}

\subsection{Nivel 3 -- Componentes}

\textbf{Pendiente para la Entrega Final, por decisión explícita tomada
durante el desarrollo de esta entrega} -- no es un olvido. El propio
\texttt{workspace.dsl} documenta el motivo: el detalle fino de
componentes del backend de Préstamos podía cambiar mientras el módulo
seguía en desarrollo activo durante esta entrega, y construir el
diagrama de Nivel~3 dos veces (una ahora, otra cuando el módulo se
estabilizara) habría sido trabajo perdido. Con el módulo de
Préstamos/Reservaciones/Multas ya completo y commiteado al cierre de
esta entrega, el Nivel~3 puede construirse sobre una base estable en la
Entrega Final.

\par\vspace{0.3cm}\noindent\colorbox{grisclaro}{\parbox{0.94\linewidth}{\vspace{0.15cm}
\textbf{Sobre la exportación a PNG.} Los diagramas C4 de Nivel~1 y
Nivel~2 existen y están validados sintácticamente
(\texttt{structurizr/cli validate}, exit code~0), pero \textbf{no se
incluye una imagen PNG exportada en este informe}: la exportación a PNG
está documentada como pendiente de integración al pipeline de CI
(Bloque~B) directamente en los comentarios de cabecera de
\texttt{docs/arquitectura/workspace.dsl}, que además deja documentado el
procedimiento completo de dos pasos (Structurizr CLI~$\to$~PlantUML~$\to$~PNG,
o revisión interactiva vía \texttt{structurizr/lite}) para cuando se
integre. Los diagramas están disponibles en formato Structurizr DSL en
\texttt{docs/arquitectura/workspace.dsl}, ejecutable/visualizable por
cualquier evaluador con Docker.
\vspace{0.15cm}}}\vspace{0.3cm}\par

\subsection{Atributos de calidad (ISO/IEC 25010)}

\texttt{docs/arquitectura/ISO25010.md} prioriza las 8 características de
calidad de producto de la norma para el contexto específico de este
proyecto (biblioteca universitaria de bajo volumen, no un sistema de alta
concurrencia). De las 8, cinco se priorizan como \textbf{Alta}
(Adecuación funcional, Usabilidad, Fiabilidad, Seguridad,
Mantenibilidad, Portabilidad -- seis, en realidad, ver el documento
fuente para el detalle exacto), con su estrategia concreta de mitigación
citada junto al ADR correspondiente cuando existe uno. Un ejemplo
representativo: \emph{Fiabilidad} se documenta honestamente como
\textbf{parcialmente atendida} -- ADR-003 deja anotado que, si Redis cae,
\texttt{JwtAuthFilter} queda sin forma de verificar revocaciones, y la
decisión \emph{fail-open} vs.\ \emph{fail-closed} sigue \textbf{pendiente
para producción}, no resuelta.

% ============================================================
% 4. MATRIZ DE TRAZABILIDAD
% ============================================================
\section{Matriz de trazabilidad}
\label{sec:matriz}

La matriz completa vive en \texttt{docs/trazabilidad/matriz.csv}
(30 filas, columnas: \texttt{id\_requisito}, \texttt{tipo},
\texttt{prioridad\_moscow}, \texttt{historia\_usuario},
\texttt{caso\_de\_uso}, \texttt{modulo\_codigo}, \texttt{endpoint\_api},
\texttt{prueba\_automatizada}, \texttt{tipo\_acceso},
\texttt{evidencia\_empirica}, \texttt{estado}) y se valida
automáticamente en cada corrida de CI vía
\texttt{scripts/validate-traceability.sh}, que rechaza el build si algún
requisito \texttt{Must} carece de historia de usuario o caso de uso, o si
algún requisito marcado \texttt{verificado} carece de prueba automatizada
real. Este informe no reproduce las 30 filas -- resume su estado
agregado.

\subsection{Resumen agregado}

\begin{center}
\begin{tabular}{|l|c|c|}
\hline
\rowcolor{verdeuteq}
\color{white}\textbf{Dimensión} & \color{white}\textbf{Categoría} & \color{white}\textbf{Cantidad} \\
\hline
\multirow{2}{*}{Tipo} & Funcional & 16 \\
 & No funcional & 14 \\
\hline
\multirow{2}{*}{Prioridad MoSCoW} & Must & 23 \\
 & Should & 7 \\
\hline
\multirow{3}{*}{Estado} & verificado & 21 \\
 & implementado & 7 \\
 & pendiente & 2 \\
\hline
\textbf{Total} & & \textbf{30} \\
\hline
\end{tabular}
\end{center}

\subsection{Cruce prioridad $\times$ estado}

\begin{center}
\begin{tabular}{|l|c|c|c|}
\hline
\rowcolor{verdeuteq}
\color{white}\textbf{Prioridad} & \color{white}\estadoVerificado & \color{white}\estadoImplementado & \color{white}\estadoPend \\
\hline
Must (23) & 21 & 2 & 0 \\
Should (7) & 0 & 5 & 2 \\
\hline
\end{tabular}
\end{center}

\textbf{Ningún requisito \texttt{Must} quedó \texttt{pendiente}}: los 23
requisitos de prioridad más alta están al menos \texttt{implementado}
(21 con evidencia empírica completa + prueba automatizada real, 2 con el
código funcionando pero evidencia empírica parcial). Los únicos 2
requisitos genuinamente \texttt{pendiente} son de prioridad
\texttt{Should}: \texttt{REQ-NF-012} (transporte TLS, mapeado a OWASP
A02) y \texttt{REQ-NF-014} (cabeceras de seguridad/CSP en
\texttt{SecurityConfig} y \texttt{nginx.conf}, mapeado a OWASP A05) --
ambos reclasificados explícitamente de \texttt{Must} a \texttt{Should}
durante esta entrega, con el criterio de que un entorno de desarrollo
sin exposición pública no justifica exigir TLS real hoy (la guía exige
un entorno públicamente accesible con TLS recién en la Entrega Final).
Este mapeo es consistente con el detalle completo del
\autoref{sec:resultados}, subsección de Seguridad.

\subsection{Trazabilidad de los dos fixes de seguridad de esta entrega}

\texttt{REQ-F-005} y \texttt{REQ-F-006} (catálogo de libros) pasaron de
\texttt{implementado} a \texttt{verificado} dentro de esta misma entrega,
como consecuencia directa del fix de rol \texttt{ADMIN} en
\texttt{LibroController} (ver \autoref{sec:resultados}): antes de ese
fix, la evidencia empírica de esos dos requisitos era parcial porque el
control de acceso real tenía un hueco no detectado hasta que se verificó
en vivo con una cuenta exclusivamente \texttt{ADMIN}.

% ============================================================
% 5. PROTOCOLO EXPERIMENTAL
% ============================================================
\section{Protocolo experimental}
\label{sec:protocolo}

\subsection{Determinismo y reproducibilidad}

Toda evidencia empírica de este proyecto sigue una convención fija,
documentada en \texttt{docs/mediciones/README.md}:

\begin{itemize}[leftmargin=1.2cm]
  \item \textbf{Cabecera obligatoria generada por script, no a mano.}
    \texttt{scripts/mediciones-header.sh} captura, en el momento exacto
    de la corrida: fecha ISO~8601~UTC real, hash corto del commit,
    y versiones exactas de Docker, Docker Compose, Java, Maven,
    PostgreSQL y Redis. Ningún archivo de evidencia completa esta
    cabecera a mano ni copia la de otro archivo.
  \item \textbf{Convención de nombres.} Evidencia puntual:
    \texttt{YYYY-MM-DD-descripcion-corta.md}. Corridas de k6:
    \texttt{kNN-runM.\{json,md\}}, donde \texttt{NN} identifica el
    escenario y \texttt{M} el número de repetición -- nunca se
    sobrescribe una corrida anterior con el mismo nombre, para poder
    comparar variabilidad entre corridas.
  \item \textbf{Semilla fija obligatoria} para cualquier dato aleatorio o
    simulado, documentada en el propio script que la usa (no solo en el
    archivo de resultados) -- convención ya aplicada preventivamente en
    \texttt{k6/opts.js} aunque las corridas de carga aún no se hayan
    ejecutado.
  \item \textbf{Archivos crudos inmutables.} La edición manual de un
    archivo de evidencia generado por una herramienta invalida esa
    evidencia -- los archivos de \texttt{docs/mediciones/jacoco/}
    (\texttt{report.xml}, \texttt{report.csv}, \texttt{html/}) son la
    salida exacta de \texttt{jacoco-maven-plugin}, sin ningún ajuste
    posterior.
\end{itemize}

\subsection{Estado de preparación por sub-bloque del Bloque C}

\begin{center}
\small
\begin{tabularx}{\linewidth}{|l|X|c|}
\hline
\rowcolor{verdeuteq}
\color{white}\textbf{Sub-bloque} & \color{white}\textbf{Configuración lista} & \color{white}\textbf{Ejecutado} \\
\hline
C.1 Rendimiento (k6) & \texttt{k6/opts.js}: perfil de 50~VUs, 30\,s de duración sostenida, con \emph{ramp-up} declarado (no salto directo a 50~VUs). 3 corridas planeadas por escenario. & No \\
\hline
C.2 Seguridad (OWASP) & 6 controles auditados en vivo + 2 correcciones con evidencia antes/después. & \textbf{Sí} \\
\hline
C.3 Usabilidad (SUS) & Plantilla de consentimiento informado (\texttt{docs/etica/consentimientos/plantilla.md}), tarea de \emph{onboarding} de 2 pasos definida, código de participante anónimo (\texttt{P01}, \texttt{P02}, \ldots). Objetivo: 10 participantes. & No \\
\hline
C.4 Cobertura (JaCoCo) & \texttt{jacoco-maven-plugin} configurado en \texttt{backend-springboot/pom.xml}, exclusiones documentadas. & \textbf{Sí} \\
\hline
C.5 Accesibilidad (Lighthouse) & No hay configuración/script propio en el repositorio todavía. & No \\
\hline
\end{tabularx}
\end{center}

\subsection{Por qué C.1, C.3 y C.5 no se ejecutaron en esta entrega}

Declarado con la misma honestidad exigida por la guía: el tiempo
disponible del equipo para esta Tercera Entrega se concentró en cerrar
primero los bloques con mayor peso de evaluación y mayor riesgo si
quedaban con evidencia fabricada o inflada -- seguridad (C.2) y cobertura
(C.4), ambos con hallazgos reales que corregir, no solo medir. Ejecutar
k6/SUS/Lighthouse \emph{de forma apresurada} solo para poder marcar la
casilla habría producido números sin valor real (una corrida de 50~VUs
contra un entorno de desarrollo sin caracterizar previamente, o una
prueba SUS con menos de los 10 participantes objetivo, no aportan
evidencia confiable) -- se prefirió declarar estos tres sub-bloques como
\textbf{no ejecutados}, con su protocolo completo y reproducible ya
preparado, en vez de generar datos de relleno. Ver
\autoref{sec:amenazas} para el análisis crítico de esta decisión.

% ============================================================
% 6. RESULTADOS POR BLOQUE
% ============================================================
\section{Resultados por bloque}
\label{sec:resultados}

\subsection{Seguridad (OWASP Top 10) -- Bloque C.2}

Se auditaron 6 controles del OWASP Top~10:2021 en vivo contra el stack
Docker real (\texttt{docker compose up -d --build}), documentados en
\texttt{docs/mediciones/sec/}. Dos de los seis hallazgos iniciales eran
gaps reales que \textbf{se corrigieron dentro de esta misma entrega},
con evidencia de antes y después contra el mismo comando de verificación.

\begin{center}
\small
\begin{tabularx}{\linewidth}{|l|X|c|}
\hline
\rowcolor{verdeuteq}
\color{white}\textbf{Control} & \color{white}\textbf{Hallazgo} & \color{white}\textbf{Estado} \\
\hline
A01 -- Control de acceso roto &
Acceso cruzado entre usuarios LECTOR (\texttt{GET~/api/v1/prestamos/usuario/\{id\}} de otro usuario) rechazado con 403 real. &
\estadoPasa \\
\hline
A02 -- Fallos criptográficos &
Sin TLS activo en el entorno de desarrollo (HTTP plano, \texttt{docker-compose.yml} sin proxy TLS) -- exigido recién en Entrega Final. Código sí declara \texttt{Secure}/\texttt{HttpOnly}/\texttt{SameSite=Strict} en la cookie del \texttt{refreshToken}, verificado en la respuesta real. &
\estadoGap \\
\hline
A03 -- Inyección &
Payloads \texttt{' OR '1'='1} y \texttt{'; DROP TABLE usuarios; --} en campos de texto libre de \texttt{/api/auth/registro}: se guardan literalmente como texto, la tabla \texttt{usuarios} permanece intacta (verificado con un login posterior). Todo el acceso a datos usa JPA parametrizado o SPs con parámetros nombrados. &
\estadoPasa \\
\hline
A05 -- Configuración de seguridad &
\texttt{X-Content-Type-Options}/\texttt{X-Frame-Options} presentes en backend y frontend. \texttt{Content-Security-Policy} ausente en ambos -- omisión de configuración real, independiente del gap de TLS. &
\estadoGap \\
\hline
A07 -- Fallos de identificación y autenticación &
Auditoría original: 6/6 intentos fallidos de login devolvían 401 sin ningún límite (sin \texttt{bucket4j} ni filtro de rate limiting). \textbf{Corregido}: \texttt{LoginRateLimiter} (Redis, clave \texttt{correo+IP}, 5 intentos / 900\,s) -- verificado en vivo: intentos 1--5 devuelven 401, el 6.\textsuperscript{o} devuelve 429 con tiempo de espera informado; reseteo confirmado tras un login exitoso. &
\estadoCorregido \\
\hline
A09 -- Fallos de registro y monitoreo &
Auditoría original: cero líneas de log y cero filas de auditoría en login exitoso o fallido (\texttt{grep} de \texttt{Logger}/\texttt{log.} en el flujo de autenticación: sin resultados). \textbf{Corregido}: logging de aplicación (\texttt{AuthService}) + \texttt{bitacora\_auditoria} (\texttt{LOGIN\_OK}/\texttt{LOGIN\_FAIL}/\texttt{LOGOUT}) con IP, timestamp y \texttt{sub}/correo -- verificado en vivo contra logs reales y la tabla real. &
\estadoCorregido \\
\hline
\end{tabularx}
\end{center}

\subsubsection*{Evidencia cruda seleccionada}

Los archivos completos de \texttt{docs/mediciones/sec/} incluyen cabeceras
HTTP íntegras, comandos exactos y análisis de cada corrida; aquí se citan
extractos representativos sin editar, para que este informe no dependa
únicamente de la tabla-resumen.

\textbf{A01 -- acceso cruzado entre usuarios LECTOR, rechazado:}
\begin{quote}\small
\begin{verbatim}
$ curl --include -s -X GET "http://localhost:8080/api/v1/prestamos/usuario/4" \
    -H "Authorization: Bearer $TOKEN_A"

HTTP/1.1 403
Content-Type: application/problem+json

{"detail":"No tiene permisos para realizar esta acción.",
 "instance":"/api/v1/prestamos/usuario/4","status":403,"title":"Forbidden"}
\end{verbatim}
\end{quote}

\textbf{A03 -- payload \texttt{'; DROP TABLE usuarios; --} en campo de
texto libre, guardado literalmente, tabla intacta:}
\begin{quote}\small
\begin{verbatim}
$ curl --include -s -X POST http://localhost:8080/api/auth/registro \
    -d "{\"nombre\":\"' OR '1'='1\",
         \"apellido\":\"'; DROP TABLE usuarios; --\", ...}"

HTTP/1.1 201
{"id":5,"nombre":"' OR '1'='1",
 "correo":"usuarioC.owasp@sgb-saas.local","roles":["LECTOR"]}

# login posterior de un usuario preexistente, para confirmar integridad:
HTTP_STATUS:200   <-- la tabla "usuarios" sigue intacta
\end{verbatim}
\end{quote}

\textbf{A07 (corregido) -- sexto intento consecutivo de login, antes 401,
ahora 429; contador real en Redis:}
\begin{quote}\small
\begin{verbatim}
--- intento 6 de 6 ---
HTTP/1.1 429
{"detail":"Demasiados intentos fallidos. Intente nuevamente en 898 segundos.",
 "instance":"/api/auth/login","status":429,"title":"Too Many Requests"}

$ redis-cli GET "login-attempts:usuarioA07fix.owasp@sgb-saas.local:172.18.0.1"
5
$ redis-cli TTL "login-attempts:usuarioA07fix.owasp@sgb-saas.local:172.18.0.1"
884
\end{verbatim}
\end{quote}

\textbf{A09 (corregido) -- login/logout ahora quedan en el log de
aplicación y en \texttt{bitacora\_auditoria}:}
\begin{quote}\small
\begin{verbatim}
2026-07-30T20:45:02.958Z INFO ... AuthService : Login exitoso:
    sub=13 correo=usuarioReseteo.owasp@sgb-saas.local ip=172.18.0.1
2026-07-30T20:45:32.583Z INFO ... AuthService : Logout:
    correo=usuarioReseteo.owasp@sgb-saas.local
    jti=649f631d-406f-4134-b7f2-4bb43abf09e2 ip=172.18.0.1

 id | usuario_id | tipo_operacion |            detalles            | ip_origen
----+------------+----------------+---------------------------------+-----------
 10 |         13 | LOGIN_OK       | Login exitoso para correo: ...  | 172.18.0.1
\end{verbatim}
\end{quote}

\subsubsection*{Detalle: A02, el gap declarado en vez de simulado}

Vale la pena destacar el tratamiento de A02 como ejemplo del criterio de
honestidad aplicado en todo el proyecto: el archivo de evidencia
correspondiente abre explícitamente diciendo que \emph{``no hay forma
honesta de verificar TLS en este entorno porque no existe TLS
activo''}, y que \emph{``cualquier evidencia que mostrara un candado o
un \texttt{https://} funcionando aquí sería fabricada''}. En vez de
omitir la sección o fabricar una captura, se documentó qué SÍ es
verificable hoy (el código ya declara el atributo \texttt{Secure}, listo
para activarse el día que exista TLS real sin cambios de código) y qué
NO (el propio TLS), dejando el segundo como gap conocido y explícitamente
fuera de alcance hasta que exista un entorno públicamente accesible.

\subsubsection*{Detalle: los dos fixes reales de esta entrega}

Más allá de la auditoría OWASP, esta entrega corrigió dos defectos reales
adicionales, detectados durante el propio proceso de verificación (no
reportados por un tercero):

\begin{itemize}[leftmargin=1.2cm]
  \item \textbf{Rol \texttt{ADMIN} faltante en \texttt{LibroController}.}
    Los 5 endpoints de creación/edición/eliminación de libros no incluían
    \texttt{ADMIN} en su lista de roles permitidos, pese a que
    \texttt{ADMIN} es el rol de mayor privilegio en la jerarquía del
    sistema (\texttt{ADMIN > GERENTE > BIBLIOTECARIO > LECTOR}). Corregido
    y verificado en vivo con una cuenta exclusivamente \texttt{ADMIN}
    (\texttt{docs/mediciones/sec/2026-07-30-owasp-a01-fix-rol-admin-libros.md}).
    Se documentó explícitamente que \texttt{PrestamoController} y
    \texttt{ReservacionController} tienen el mismo hueco y quedan
    \textbf{sin corregir}, fuera del alcance de ese fix puntual.
  \item \textbf{\texttt{POST /api/auth/refresh} respondía 500 en vez de
    401} cuando el \texttt{refreshToken} de la cookie era
    inválido/expirado: una \texttt{RuntimeException} genérica sin
    \texttt{@ExceptionHandler} dedicado caía en el handler de último
    recurso. Corregido con una excepción específica
    (\texttt{RefreshTokenInvalidoException}) y su handler en
    \texttt{GlobalExceptionHandler}, devolviendo 401 vía
    \texttt{ProblemDetail} (RFC~7807).
\end{itemize}

\subsection{Cobertura de código (JaCoCo) -- Bloque C.4}

Reporte real generado por \texttt{jacoco-maven-plugin} (versión 0.8.12)
sobre el dominio/servicios del backend, excluyendo \texttt{dto/}
(records planos sin lógica), \texttt{entity/} (entidades JPA con
\texttt{@Data}, sin métodos de negocio propios) y \texttt{config/}
(definición de beans de framework) -- criterio documentado también como
comentario en \texttt{backend-springboot/pom.xml}. Artefactos completos
versionados en \texttt{docs/mediciones/jacoco/} (\texttt{report.xml},
\texttt{report.csv}, \texttt{html/}).

\subsubsection*{Totales agregados}

\begin{center}
\begin{tabular}{|l|c|c|c|}
\hline
\rowcolor{verdeuteq}
\color{white}\textbf{Métrica} & \color{white}\textbf{Cubierto} & \color{white}\textbf{Total} & \color{white}\textbf{\%} \\
\hline
Líneas & 380 & 501 & \textbf{75{,}85\,\%} \\
Ramas & 50 & 102 & 49{,}02\,\% \\
Complejidad & 111 & 190 & 58{,}42\,\% \\
\hline
\end{tabular}
\end{center}

El objetivo de esta entrega (\textbf{$\geq$60\,\% de líneas}) se supera
con margen, y el resultado ya alcanza el umbral fijado para la Entrega
Final ($\geq$70\,\%). Verificado con \texttt{./mvnw clean verify}:
\texttt{BUILD SUCCESS}, 79/79 tests (0 fallos, 0 errores).

\subsubsection*{Desglose por paquete}

\begin{center}
\begin{tabular}{|l|c|c|}
\hline
\rowcolor{verdeuteq}
\color{white}\textbf{Paquete} & \color{white}\textbf{Líneas} & \color{white}\textbf{Ramas} \\
\hline
\texttt{security} & 86/90 (95{,}6\,\%) & 15/22 (68{,}2\,\%) \\
\texttt{scheduling} & 7/7 (100{,}0\,\%) & -- \\
\texttt{controller} & 46/66 (69{,}7\,\%) & 3/4 (75{,}0\,\%) \\
\texttt{service} & 224/304 (73{,}7\,\%) & 30/62 (48{,}4\,\%) \\
\texttt{exception} & 17/34 (50{,}0\,\%) & 2/14 (14{,}3\,\%) \\
\hline
\end{tabular}
\end{center}

\subsubsection*{Historia: de 53{,}69\,\% a 75{,}85\,\%}

La primera corrida real de JaCoCo en esta entrega arrojó 53{,}69\,\% de
líneas -- por debajo del objetivo. En vez de agregar tests triviales
para inflar el número (instrucción explícita seguida durante todo el
proceso), se identificaron 5 clases de autenticación/autorización con
cobertura casi nula \emph{dentro del proceso de test} porque los tests
existentes las mockeaban en vez de ejercitarlas: \texttt{JwtService},
\texttt{JwtAuthFilter}, \texttt{UserDetailsServiceImpl},
\texttt{AuthController} y \texttt{GlobalExceptionHandler}. Se agregaron
4 clases de test nuevas que ejercitan la implementación \emph{real} de
cada una (\texttt{JwtServiceTest} unitario puro sin contexto Spring,
\texttt{JwtAuthFilterTest} con el filtro real y solo Redis/BD mockeados,
\texttt{UserDetailsServiceImplTest}, y \texttt{AuthControllerTest} vía
\texttt{MockMvc} + \texttt{springSecurity()} real, mismo patrón que
\texttt{LibroControllerSecurityTest}):

\begin{center}
\begin{tabular}{|l|c|c|}
\hline
\rowcolor{verdeuteq}
\color{white}\textbf{Clase} & \color{white}\textbf{Antes} & \color{white}\textbf{Después} \\
\hline
\texttt{JwtService} & 4{,}8\,\% & \textbf{95{,}2\,\%} \\
\texttt{JwtAuthFilter} & 18{,}2\,\% & \textbf{100{,}0\,\%} \\
\texttt{UserDetailsServiceImpl} & 7{,}1\,\% & \textbf{100{,}0\,\%} \\
\texttt{AuthController} & 0{,}0\,\% & \textbf{100{,}0\,\%} \\
\texttt{GlobalExceptionHandler} & 8{,}8\,\% & \textbf{50{,}0\,\%} \\
\hline
\end{tabular}
\end{center}

\texttt{GlobalExceptionHandler} se dejó deliberadamente en 50\,\%: las
líneas restantes corresponden a \texttt{handleNotFound}
(\texttt{EntityNotFoundException}, no aplica a ningún flujo de
\texttt{AuthController}) y al \emph{catch-all} de último recurso
(\texttt{handleGenerica}) -- forzar esas líneas habría exigido tests
artificiales fuera del alcance de auth/seguridad, exactamente el tipo de
relleno que se pidió evitar.

\subsubsection*{Detalle completo por clase (las 20 clases analizadas)}

Tabla completa de \texttt{docs/mediciones/jacoco/report.csv}, ordenada de
menor a mayor cobertura de líneas -- ninguna clase se omite:

\begin{center}
\small
\begin{longtable}{|l|l|c|}
\hline
\rowcolor{verdeuteq}
\color{white}\textbf{Paquete} & \color{white}\textbf{Clase} & \color{white}\textbf{\% líneas} \\
\hline
\endhead
\texttt{controller} & \texttt{PrestamoController} & 25{,}0\,\% \\
\hline
\texttt{controller} & \texttt{MultaController} & 42{,}9\,\% \\
\hline
\texttt{controller} & \texttt{ReservacionController} & 42{,}9\,\% \\
\hline
\texttt{controller} & \texttt{TestController} & 50{,}0\,\% \\
\hline
\texttt{exception} & \texttt{GlobalExceptionHandler} & 50{,}0\,\% \\
\hline
\texttt{service} & \texttt{PrestamoService} & 66{,}1\,\% \\
\hline
\texttt{service} & \texttt{AuthService} & 67{,}5\,\% \\
\hline
\texttt{service} & \texttt{MultaService} & 72{,}7\,\% \\
\hline
\texttt{service} & \texttt{LibroService} & 75{,}7\,\% \\
\hline
\texttt{controller} & \texttt{LibroController} & 80{,}0\,\% \\
\hline
\texttt{security} & \texttt{LoginRateLimiter} & 83{,}3\,\% \\
\hline
\texttt{service} & \texttt{ReservacionService} & 90{,}2\,\% \\
\hline
\texttt{security} & \texttt{JwtService} & 95{,}2\,\% \\
\hline
\texttt{controller} & \texttt{AuthController} & 100{,}0\,\% \\
\hline
\texttt{scheduling} & \texttt{ReservacionScheduler} & 100{,}0\,\% \\
\hline
\texttt{security} & \texttt{JwtAuthFilter} & 100{,}0\,\% \\
\hline
\texttt{security} & \texttt{UserDetailsServiceImpl} & 100{,}0\,\% \\
\hline
\texttt{service} & \texttt{CorreoYaRegistradoException} & 100{,}0\,\% \\
\hline
\texttt{service} & \texttt{LoginRateLimitExcedidoException} & 100{,}0\,\% \\
\hline
\texttt{service} & \texttt{RefreshTokenInvalidoException} & 100{,}0\,\% \\
\hline
\end{longtable}
\end{center}

\subsection{Rendimiento -- Bloque C.1 (k6)}

\estadoPendiente

\textbf{No ejecutado en esta entrega.} La configuración base ya está
lista en \texttt{k6/opts.js}: perfil de 50~VUs, 30\,s de duración
sostenida, con etapas de \emph{ramp-up}/\emph{ramp-down} declaradas en
vez de saltar directo a la carga máxima, y 3 corridas planeadas por
escenario para poder comparar variabilidad
(\texttt{docs/mediciones/perf/kNN-runM.json}). El script de prueba en sí
(las funciones que llaman a los endpoints, los \emph{checks} y los
\emph{thresholds} de negocio) queda pendiente de escribir antes de la
primera corrida real. Ningún número de latencia/throughput se reporta en
este informe para no fabricar evidencia que no existe.

\subsection{Usabilidad -- Bloque C.3 (SUS)}

\estadoPendiente

\textbf{No ejecutado en esta entrega.} El protocolo completo está
preparado: plantilla de consentimiento informado en
\texttt{docs/etica/consentimientos/plantilla.md} (código de participante
anónimo tipo \texttt{P01}, tarea de \emph{onboarding} de dos pasos --
iniciar sesión con cuenta de prueba, registrarse como usuario nuevo --
sin ayuda del equipo salvo bloqueo total), objetivo de 10 participantes,
y la política de anonimización ya declarada en
\texttt{docs/etica/ETHICS.md} (formularios firmados nunca se committean
al repositorio; solo la plantilla en blanco). Ningún puntaje SUS se
reporta en este informe.

\subsection{Accesibilidad -- Bloque C.5 (Lighthouse)}

\estadoPendiente

\textbf{No ejecutado en esta entrega.} A diferencia de C.1 y C.3, este
sub-bloque todavía no tiene ni siquiera configuración base preparada en
el repositorio -- es el sub-bloque con menor avance de los tres
pendientes. Queda como trabajo a iniciar desde cero para la Entrega
Final: definir el conjunto de páginas del frontend Angular a auditar y
los umbrales mínimos aceptables por categoría (Accesibilidad,
Rendimiento, Buenas prácticas, SEO).

% ============================================================
% 7. AMENAZAS A LA VALIDEZ
% ============================================================
\section{Amenazas a la validez}
\label{sec:amenazas}

Análisis crítico real de las limitaciones de esta entrega, no una
sección de relleno:

\begin{enumerate}[leftmargin=1.2cm]
  \item \textbf{Tres de cinco sub-bloques de evidencia empírica del
    Bloque~C sin ejecutar.} Rendimiento (C.1), usabilidad (C.3) y
    accesibilidad (C.5) quedaron como protocolo/configuración lista pero
    sin corrida real, por restricción de tiempo del equipo dentro de esta
    entrega (ver \autoref{sec:protocolo}). Esto significa que las
    afirmaciones de \emph{eficiencia de desempeño} en
    \texttt{docs/arquitectura/ISO25010.md} se apoyan únicamente en una
    medición puntual de caché (\texttt{~170ms} \emph{miss} vs.\
    \texttt{~30ms} \emph{hit}), no en una prueba de carga formal -- no
    hay evidencia todavía de cómo se comporta el sistema bajo 50~VUs
    concurrentes.
  \item \textbf{Escala del equipo limita la escala de las pruebas de
    usabilidad.} El equipo tiene 3 integrantes, todos con roles activos
    de desarrollo durante esta entrega -- reclutar y coordinar 10
    participantes externos para SUS (objetivo declarado en el protocolo)
    compite directamente por el mismo tiempo limitado que exige cerrar
    los demás bloques. Esta es una limitación estructural del proyecto,
    no solo de esta entrega puntual.
  \item \textbf{Auditoría de seguridad manual, no exhaustiva.} Las 6
    verificaciones OWASP se hicieron con \texttt{curl} dirigido a
    escenarios específicos elegidos por el equipo, no con una herramienta
    de \emph{fuzzing}/escaneo automatizado (ej.\ OWASP ZAP) que explore
    superficie de ataque no anticipada. Además, varios hallazgos se
    verificaron solo en un endpoint representativo y se documentó
    explícitamente que el mismo patrón se asume (no se re-verificó) en
    endpoints hermanos -- ej.\ A01 se probó solo contra
    \texttt{PrestamoController}, asumiendo que
    \texttt{MultaController}/\texttt{ReservacionController} comparten el
    mecanismo (\texttt{validarAccesoUsuario}), sin repetir la prueba en
    cada uno.
  \item \textbf{Gap de validación en \texttt{ReservacionService}
    (hallazgo nuevo, no corregido).} Detectado al construir la colección
    Postman de esta entrega:
    \texttt{ReservacionService.crear()} no valida que \texttt{libroId}
    exista antes de guardar (a diferencia de \texttt{PrestamoController},
    que sí pasa por un procedimiento almacenado con validación
    \texttt{LB404}). Una reservación con un libro inexistente no
    responde 404: la restricción de llave foránea de PostgreSQL revienta
    como excepción no mapeada, y \texttt{GlobalExceptionHandler} la
    convierte en \textbf{500 Internal Server Error}. Verificado en vivo
    y documentado en \texttt{docs/postman/coleccion.json}, no corregido
    -- fuera del alcance del prompt que lo detectó.
  \item \textbf{\texttt{CONTRIBUTORS.md} desactualizado respecto al
    estado real del repositorio.} El propio archivo declara, en su nota
    de cabecera, que el backend de Préstamos/Reservas y los ajustes de
    frontend correspondientes aún no estaban commiteados al momento de
    escribirse -- al cierre de esta entrega, ese trabajo sí está
    commiteado, pero el archivo no se ha vuelto a revisar. Se documenta
    aquí en vez de inventar una actualización que no se verificó línea
    por línea contra \texttt{git log} en el momento de escribir este
    informe.
  \item \textbf{TLS y CSP diferidos, no resueltos.} OWASP A02 (TLS en
    tránsito) y parte de A05 (\emph{Content-Security-Policy}) quedan como
    gap conocido explícito -- correcto para un entorno de desarrollo sin
    exposición pública, pero significa que la postura de seguridad
    completa del sistema no puede evaluarse hasta que exista el entorno
    públicamente accesible que exige la Entrega Final.
\end{enumerate}

% ============================================================
% 8. CONTRIBUCIONES Y ÉTICA
% ============================================================
\section{Contribuciones y declaración ética}

\subsection{Contribuciones (taxonomía CRediT)}

El detalle completo de roles por integrante, basado en trabajo
realmente commiteado (no en el rol nominal de equipo), vive en
\texttt{CONTRIBUTORS.md} siguiendo la taxonomía CRediT (14 roles
estándar, \url{https://credit.niso.org/}). Resumen:

\begin{itemize}[leftmargin=1.2cm]
  \item \textbf{Marlon Loor Medranda} (Tech Lead / DevOps / Seguridad):
    \emph{Software} (JWT+RBAC, cookies \texttt{HttpOnly}, blacklist de
    tokens, caché con TTL externo, \texttt{GlobalExceptionHandler}),
    \emph{Project administration}, \emph{Supervision} (revisión de
    hallazgos de QA sobre trabajo de otros integrantes) y
    \emph{Validation} (verificación en vivo de cada cambio de
    seguridad/caché contra el stack Docker real).
  \item \textbf{Irvin Cajas Ibarra} (Backend): \emph{Software} -- CRUD de
    libros y el módulo de Préstamos/Reservas/Multas.
  \item \textbf{Moises Panama Murillo} (Frontend): \emph{Software} --
    módulos Angular de autenticación (login/registro, guards,
    interceptor JWT) y CRUD de libros en el frontend.
\end{itemize}

\subsection{Declaración ética}

\texttt{docs/etica/ETHICS.md} cubre las cuatro declaraciones exigidas por
el Bloque~F de la guía:

\begin{enumerate}[leftmargin=1.2cm]
  \item \textbf{Procedencia de datos.} Los 5 libros de ejemplo de
    \texttt{db/seed.sql} \textbf{no son ficticios}: son libros reales
    publicados, con ISBN real (ej.\ \emph{Clean Code}, ISBN
    9780132350884). Los metadatos bibliográficos (ISBN, título, autor,
    editorial) son hechos públicos de catalogación, transcritos
    manualmente por el equipo -- sin \emph{scraping} ni consumo de API de
    terceros con licencia. El resumen de una línea por libro es redacción
    original del equipo. El contenido íntegro de ninguna obra está
    incluido en el repositorio.
  \item \textbf{Tratamiento de datos personales.} Contraseñas hasheadas
    con BCrypt costo~12, nunca en texto plano. Auditoría explícita del
    código fuente confirmó que \texttt{password\_hash} no se expone por
    ningún endpoint ni DTO (\texttt{@JsonIgnore} sobre el campo, y el
    único DTO de registro no incluye ese campo en su definición).
  \item \textbf{Consentimiento informado para SUS.} Mecanismo declarado
    \emph{antes} de la primera prueba, no retroactivamente -- ver
    \autoref{sec:protocolo}.
  \item \textbf{Ausencia de datos identificables en el repositorio
    público.} \texttt{docs/mediciones/} solo contiene evidencia técnica
    (códigos de estado HTTP, \emph{timings}, TTLs de Redis); los
    consentimientos firmados de participantes reales de SUS nunca se
    committean, solo la plantilla en blanco.
\end{enumerate}

% ============================================================
% 9. REFERENCIAS
% ============================================================
\section{Referencias}

\begin{enumerate}[leftmargin=1.2cm]
  \item ISO/IEC/IEEE 29148:2018 -- \emph{Systems and software
    engineering -- Life cycle processes -- Requirements engineering}.
  \item ISO/IEC 25010:2011 -- \emph{Systems and software Quality
    Requirements and Evaluation (SQuaRE) -- Product quality model}.
  \item RFC~7519 -- \emph{JSON Web Token (JWT)}, M.\ Jones, J.\
    Bradley, N.\ Sakimura, IETF, mayo 2015.
  \item RFC~7807 -- \emph{Problem Details for HTTP APIs}, M.\ Nottingham,
    E.\ Wilde, IETF, marzo 2016.
  \item OWASP Top~10:2021 -- \emph{The Ten Most Critical Web Application
    Security Risks}, OWASP Foundation.
    \url{https://owasp.org/Top10/}
  \item Semantic Versioning 2.0.0. \url{https://semver.org/}
  \item Keep a Changelog 1.1.0. \url{https://keepachangelog.com/}
  \item CRediT -- Contributor Roles Taxonomy. NISO/CASRAI.
    \url{https://credit.niso.org/}
  \item Brooke, J.\ (1996). \emph{SUS: A ``quick and dirty'' usability
    scale}. En P.W.\ Jordan et al.\ (Eds.), \emph{Usability Evaluation in
    Industry}. Taylor \& Francis.
  \item Brown, S.\ -- \emph{The C4 model for visualising software
    architecture}. \url{https://c4model.com/}
  \item Google -- \emph{Lighthouse}, herramienta de auditoría automatizada
    de calidad web. \url{https://developer.chrome.com/docs/lighthouse/}
  \item Grafana Labs -- \emph{k6}, herramienta de pruebas de carga.
    \url{https://k6.io/}
\end{enumerate}

\vspace{1cm}
\noindent\rule{\linewidth}{0.4pt}

\vspace{0.3cm}
\noindent{\small Este informe se generó a partir de evidencia real
versionada en el repositorio del proyecto
(\url{https://github.com/mloorm14/sgb-saas}, tag \texttt{v0.9.0-rc},
commit \texttt{c6b372c}). Ningún dato de rendimiento, usabilidad o
accesibilidad reportado en este documento fue inventado o extrapolado --
donde la evidencia no existe (Bloques C.1, C.3, C.5), se declaró
explícitamente como no ejecutada.}

\end{document}
